
\subsection{算子 $\mathfrak r$ 的代数表述 II：分叉论证}
\label{subsec:7-4-8}

现在回到几何情形。
设 $M$ 中一对相对 spin 的拉格朗日子流形
$L^{(1)},L^{(0)}$ 干净相交。
由定理~\ref{thm:A}，
存在含幺滤过 $A_\infty$ 代数
\[
H(L^{(1)};\Lambda_{0,\mathrm{nov}}),
\qquad
H(L^{(0)};\Lambda_{0,\mathrm{nov}}),
\]
而由定理~\ref{thm:F}，
\begin{equation*}
H(L^{(1)},L^{(0)};\Lambda_{0,\mathrm{nov}})
=
\bigoplus_h
H(R_h,\Theta^-_h;\Lambda_{0,\mathrm{nov}})
[\mu(R_h)]
\end{equation*}
上具有含幺滤过 $A_\infty$ 双模结构。

此外，由定理~\ref{thm:7-4-118}，
存在严格、有隙且滤过的 $L_\infty$ 同态
\[
\begin{aligned}
\mathfrak q^\circ_{(i)}\colon
H(M;\Lambda_{0,\mathrm{nov}})[1]
&\longrightarrow
\operatorname{Der}B
\bigl(H(L^{(i)};\Lambda_{0,\mathrm{nov}})[1]\bigr),\\
\mathfrak q^{\circ+}_{(i)}\colon
H(M;\Lambda_{0,\mathrm{nov}})[1]
&\longrightarrow
\operatorname{Der}B^{\mathrm{unit}}
\bigl(H(L^{(i)};\Lambda_{0,\mathrm{nov}})[1]\bigr).
\end{aligned}
\]
这里引用定理~\ref{thm:7-4-118}
\translatornote{源书误写“定理 7.4.102”；(7.4.102) 只是定理 7.4.41 证明中的一个公式，而这里所用的 $\mathfrak q^\circ$ 及其含幺版本正由定理 7.4.118 给出。见疑误 A-0599。}。

\begin{theorem}{7.4.154}\label{thm:7-4-154}
存在严格、有隙且滤过的 $L_\infty$ 同态
\[
\mathfrak r^\circ\colon
H(M;\Lambda_{0,\mathrm{nov}})[1]
\longrightarrow
\operatorname{Der}B
\bigl(
H(L^{(1)},L^{(0)};\Lambda_{0,\mathrm{nov}})[1]
\bigr),
\]
使下列图式交换：

\par\medskip\noindent\begin{minipage}{\textwidth}
\centering
\resizebox{.90\textwidth}{!}{$\displaystyle
\begin{array}{ccc}
H(M)[1]
&\xrightarrow{\ \mathfrak r^\circ\ }&
\operatorname{Der}B
\bigl(H(L^{(1)},L^{(0)})[1]\bigr)
\\[-.1em]
\Big\Vert&&\Big\downarrow
\\[-.1em]
H(M)[1]
&\xrightarrow{\ \mathfrak q^\circ_{(1)}
\oplus\mathfrak q^\circ_{(0)}\ }&
\operatorname{Der}B\bigl(H(L^{(1)})[1]\bigr)
\oplus
\operatorname{Der}B\bigl(H(L^{(0)})[1]\bigr).
\end{array}$}
\par\smallskip
{\renewcommand{\thefigure}{7.4.5}\refstepcounter{figure}
\normalfont\bfseries 图式~\thefigure\label{diag:7-4-5}}
\end{minipage}\par\medskip

图式~\ref{diag:7-4-5} 中省略了系数环
$\Lambda_{0,\mathrm{nov}}$。
含幺版本亦成立。
\end{theorem}

\begin{proof}
这是定理~\ref{thm:3-8-71} 的重新表述。
沿用该处的 $\mathfrak r_{\ell;k_1,k_0}$，
对 $\ell>0$ 置
\[
\mathfrak r^{\circ,\mathrm{chain}}_{0,\ell}
(\mathbf y)(\mathbf x,v,\mathbf z)
=
\mathfrak r_{\ell;k_1,k_0}
\bigl(\operatorname{Op}(\mathbf y)\bigr)
(\mathbf x,v,\mathbf z).
\]
于是它与 $\mathfrak q^\circ$ 一同给出一个图式，
该图式由图式~\ref{diag:7-4-5}
把所有 $H$ 换成 $C$ 得到。
事实上，式~\eqref{eq:3-8-72}
与式 (3.8.73.1) 合起来，
等价于
$\mathfrak r^{\circ,\mathrm{chain}}_{0,\ell}$
是 $L_\infty$ 同态这一条件。
除符号外，这一点由定义相当显然；
符号可仿照命题~\ref{prop:7-4-34}
与推论~\ref{cor:7-4-40} 的证明核对。
最后应用定理~\ref{thm:7-4-149}，
从链层过渡到同调层。
\end{proof}

下面讨论定理~\ref{thm:7-4-154}
中的结构对参数选取的依赖。
如第~\ref{sec:5-3} 节已经提到的，
须考虑两种情形：
一种是用同一个辛微分同胚移动两个拉格朗日子流形，
另一种是用两个相互独立的哈密顿同痕移动它们。
本小节余下部分讨论第一种情形，
下一小节讨论第二种情形。

设
$\psi\colon M\to M'$ 是辛微分同胚，
并置
$\psi(L^{(i)})=L^{(i)\prime}$。
由定理~\ref{thm:7-4-120}，存在
\begin{equation}
\psi_*^{(i)}\colon
\operatorname{Der}B
\bigl(H(L^{(i)};\Lambda_{0,\mathrm{nov}})[1]\bigr)
\longrightarrow
\operatorname{Der}B
\bigl(H(L^{(i)\prime};\Lambda_{0,\mathrm{nov}})[1]\bigr),
\tag{7.4.155}\label{eq:7-4-155}
\end{equation}
并且它与 $\mathfrak q^\circ_{(i)}$ 交换。

\begin{theorem}{7.4.156}\label{thm:7-4-156}
\liftIndex{I203}
存在严格、有隙且滤过的 $L_\infty$ 同伦等价
\[
\begin{aligned}
\psi_*\colon
&\operatorname{Der}B
\bigl(H(L^{(1)},L^{(0)};\Lambda_{0,\mathrm{nov}})[1]\bigr)\\
&\longrightarrow
\operatorname{Der}B
\bigl(H(L^{(1)\prime},L^{(0)\prime};
\Lambda_{0,\mathrm{nov}})[1]\bigr),
\end{aligned}
\]
使图式~\ref{diag:7-4-6}、\ref{diag:7-4-7}
在同伦意义下交换。

\par\medskip\noindent\begin{minipage}{\textwidth}
\centering
\resizebox{.98\textwidth}{!}{$\displaystyle
\begin{array}{ccc}
\operatorname{Der}B
\bigl(H(L^{(1)},L^{(0)};\Lambda_{0,\mathrm{nov}})[1]\bigr)
&\xrightarrow{\ \psi_*\ }&
\operatorname{Der}B
\bigl(H(L^{(1)\prime},L^{(0)\prime};
\Lambda_{0,\mathrm{nov}})[1]\bigr)
\\[-.1em]
\Big\downarrow&&\Big\downarrow
\\[-.1em]
\begin{gathered}
\operatorname{Der}B
\bigl(H(L^{(1)};\Lambda_{0,\mathrm{nov}})[1]\bigr)\\
\oplus\\
\operatorname{Der}B
\bigl(H(L^{(0)};\Lambda_{0,\mathrm{nov}})[1]\bigr)
\end{gathered}
&\xrightarrow{\ \psi_*^{(1)}\oplus\psi_*^{(0)}\ }&
\begin{gathered}
\operatorname{Der}B
\bigl(H(L^{(1)\prime};\Lambda_{0,\mathrm{nov}})[1]\bigr)\\
\oplus\\
\operatorname{Der}B
\bigl(H(L^{(0)\prime};\Lambda_{0,\mathrm{nov}})[1]\bigr).
\end{gathered}
\end{array}$}
\par\smallskip
{\renewcommand{\thefigure}{7.4.6}\refstepcounter{figure}
\normalfont\bfseries 图式~\thefigure\label{diag:7-4-6}}
\end{minipage}\par\medskip

\par\medskip\noindent\begin{minipage}{\textwidth}
\centering
\resizebox{.93\textwidth}{!}{$\displaystyle
\begin{array}{ccc}
H(M;\Lambda_{0,\mathrm{nov}})
&\xrightarrow{\ \psi_*\ }&
H(M';\Lambda_{0,\mathrm{nov}})
\\[-.1em]
{\scriptstyle\mathfrak r^\circ}\Big\downarrow
&&
\Big\downarrow{\scriptstyle\mathfrak r^\circ}
\\[-.1em]
\operatorname{Der}B
\bigl(H(L^{(1)},L^{(0)};\Lambda_{0,\mathrm{nov}})[1]\bigr)
&\xrightarrow{\ \psi_*\ }&
\operatorname{Der}B
\bigl(H(L^{(1)\prime},L^{(0)\prime};
\Lambda_{0,\mathrm{nov}})[1]\bigr).
\end{array}$}
\par\smallskip
{\renewcommand{\thefigure}{7.4.7}\refstepcounter{figure}
\normalfont\bfseries 图式~\thefigure\label{diag:7-4-7}}
\end{minipage}\par\medskip

图式~\ref{diag:7-4-7} 左下角写作完整的
$\operatorname{Der}B$
\translatornote{源书在左下角误排为 $\operatorname{Der}B_0$，而右下角及定理 7.4.156 的源、靶均为完整的 $\operatorname{Der}B$；此外 $\mathfrak r^\circ$ 还含由两端 $\mathfrak q^\circ_{(i)}$ 给出的分量，不能取值于只含中间分量的理想 $\operatorname{Der}B_0$。见疑误 A-0600。}。
含幺版本亦成立。
\end{theorem}

\begin{proof}
这里采用的论证基于第~\ref{subsec:7-2-14} 小节
说明的分叉法。
考虑该小节中的一族相容殆复结构
$\{J_\rho\}_\rho$，
并使用注~\ref{rem:7-4-106} 中的模空间
\[
\mathcal M_{k+1,\ell}
\bigl(\{J_\rho\}_\rho,L^{(i)\prime},\beta;
\operatorname{twp}(\rho)\bigr).
\]
下面略去从 $A_{n,K}$ 过渡到 $A_\infty$ 的过程；
相关说明见注~\ref{rem:7-4-124} (1)。

仿照第~\ref{subsec:5-3-3}、\ref{subsec:7-2-10}
小节，取链的可数集
$\mathcal X([0,1]\times R_h)$，
并用它定义 $\Lambda_{0,\mathrm{nov}}$-模
\[
C_{\mathrm{para}}
(L^{(1)\prime},L^{(0)\prime};\Lambda_{0,\mathrm{nov}}).
\]
见图式~\ref{diag:5-3-3}。
假设与条件 (5.3.35.3) 类似的条件，
便有同态
\begin{equation}
\begin{aligned}
\operatorname{Eval}_{\rho=\rho_0}\colon
C_{\mathrm{para}}
(L^{(1)\prime},L^{(0)\prime};\Lambda_{0,\mathrm{nov}})
\longrightarrow
C^{\rho=\rho_0}
(L^{(1)\prime},L^{(0)\prime};\Lambda_{0,\mathrm{nov}}).
\end{aligned}
\tag{7.4.157}\label{eq:7-4-157}
\end{equation}
这里
\[
\begin{aligned}
C^{\rho=0}
(L^{(1)\prime},L^{(0)\prime};\Lambda_{0,\mathrm{nov}})
&=
\psi_*\bigl(
C(L^{(1)},L^{(0)};\Lambda_{0,\mathrm{nov}})
\bigr),
\end{aligned}
\]
而
$C(L^{(1)},L^{(0)};\Lambda_{0,\mathrm{nov}})$
正是定义图式~\ref{diag:7-4-7} 第一列中
$\mathfrak r^\circ$ 所用的链复形。
另一方面，
$C^{\rho=1}
(L^{(1)\prime},L^{(0)\prime};\Lambda_{0,\mathrm{nov}})$
是定义图式~\ref{diag:7-4-7} 第二列中
$\mathfrak r^\circ$ 所用的链复形
$C(L^{(1)\prime},L^{(0)\prime};\Lambda_{0,\mathrm{nov}})$。

现在取逐时积模空间
\begin{equation}
\begin{aligned}
\mathcal M_{k_1,k_0}
\bigl(&L^{(1)\prime},L^{(0)\prime};
[h'_1,w'_1],[h'_2,w'_2];
\{J_\rho\}_\rho,\operatorname{twp}(\rho)\bigr),
\end{aligned}
\tag{7.4.158.1}\label{eq:7-4-158-1}
\end{equation}
它定义于式~\eqref{eq:5-3-31}。
式~\eqref{eq:5-3-31} 中的族
$\{J_{\infty,t,s}\}_t$
在这里取为
$J_{\infty,t,\rho}=J_\rho$。

注意，
\[
\mathcal M_{k+1,0}
\bigl(\{J_\rho\}_\rho,L^{(i)\prime},\beta;
\operatorname{twp}(\rho)\bigr)
\]
在模空间~\eqref{eq:7-4-158-1}
的紧化中作为一个纤维积因子出现。
为~\eqref{eq:7-4-158-1} 选取扰动，
使它与上述模空间上给定的扰动（多重截面）相容。
由此在
$C_{\mathrm{para}}
(L^{(1)\prime},L^{(0)\prime};\Lambda_{0,\mathrm{nov}})$
上定义一个滤过 $A_\infty$ 双模结构，
其左、右代数分别为
\[
C([0,1]\times L^{(1)\prime};\Lambda_{0,\mathrm{nov}}),
\qquad
C([0,1]\times L^{(0)\prime};\Lambda_{0,\mathrm{nov}}).
\]
还可选择扰动，使式~\eqref{eq:7-4-157}
成为滤过 $A_\infty$ 双模同态。
于是得到下列交换图式；其中省略
$\Lambda_{0,\mathrm{nov}}$：

\par\medskip\noindent\begin{minipage}{\textwidth}
\centering
\resizebox{.94\textwidth}{!}{$\displaystyle
C(L^{(1)},L^{(0)})
\xleftarrow{\ \operatorname{Eval}_{\rho=0}\ }
C_{\mathrm{para}}(L^{(1)\prime},L^{(0)\prime})
\xrightarrow{\ \operatorname{Eval}_{\rho=1}\ }
C(L^{(1)\prime},L^{(0)\prime})$}
\par\smallskip
{\renewcommand{\thefigure}{7.4.8}\refstepcounter{figure}
\normalfont\bfseries 图式~\thefigure\label{diag:7-4-8}}
\end{minipage}\par\medskip

接着以显然的方式在~\eqref{eq:7-4-158-1}
中加入 $\ell$ 个内部标记点，得到模空间
\begin{equation}
\begin{aligned}
\mathcal M_{k_1,k_0;\ell}
\bigl(&L^{(1)\prime},L^{(0)\prime};
[h'_1,w'_1],[h'_2,w'_2];
\{J_\rho\}_\rho,\operatorname{twp}(\rho)\bigr).
\end{aligned}
\tag{7.4.158.2}\label{eq:7-4-158-2}
\end{equation}
注意，
\[
\mathcal M_{k+1,\ell}
\bigl(\{J_\rho\}_\rho,L^{(i)\prime},\beta;
\operatorname{twp}(\rho)\bigr)
\]
在模空间~\eqref{eq:7-4-158-2}
的紧化中作为一个纤维积因子出现。
定义 $\mathfrak q^\circ_{(i)}$ 时，
我们已在第~\ref{subsec:7-2-13} 小节
及式~\eqref{eq:7-4-122} 中
为这个模空间选取了扰动。
为~\eqref{eq:7-4-158-2} 选取与之相容的扰动。
仿照第~\ref{subsec:3-8-7}、\ref{subsec:3-8-8}
小节以及定理~\ref{thm:7-4-154} 的证明，
由此得到 $\mathfrak r^\circ_{\mathrm{para}}$，
使下列图式交换：

\par\medskip\noindent\begin{minipage}{\textwidth}
\centering
\resizebox{.93\textwidth}{!}{$\displaystyle
\begin{array}{ccc}
H(M)
&\xrightarrow{\ \mathfrak r^\circ\ }&
\operatorname{Der}B
\bigl(C(L^{(1)},L^{(0)})[1]\bigr)
\\[-.1em]
{\scriptstyle\psi_*}\Big\downarrow
&&
\Big\uparrow{\scriptstyle\operatorname{Eval}_{\rho=0}}
\\[-.1em]
H(M')
&\xrightarrow{\ \mathfrak r^\circ_{\mathrm{para}}\ }&
\operatorname{Der}B
\bigl(C_{\mathrm{para}}
(L^{(1)\prime},L^{(0)\prime})[1]\bigr)
\\[-.1em]
\Big\Vert
&&
\Big\downarrow{\scriptstyle\operatorname{Eval}_{\rho=1}}
\\[-.1em]
H(M')
&\xrightarrow{\ \mathfrak r^\circ\ }&
\operatorname{Der}B
\bigl(C(L^{(1)\prime},L^{(0)\prime})[1]\bigr).
\end{array}$}
\par\smallskip
{\renewcommand{\thefigure}{7.4.9}\refstepcounter{figure}
\normalfont\bfseries 图式~\thefigure\label{diag:7-4-9}}
\end{minipage}\par\medskip

最后应用定理~\ref{thm:7-4-149}，
把图式~\ref{diag:7-4-9} 中的 $C$ 换成 $H$。
定理~\ref{thm:7-4-156} 得证。
\end{proof}

\begin{remark}{7.4.159}\label{rem:7-4-159}
事实上，从 $A_{n,K}$ 陈述过渡到
$A_\infty$ 陈述时，
我们使用了定理~\ref{thm:3-8-71} 的部分证明之重新表述，
即定理~\ref{thm:7-4-154}
与定理~\ref{thm:7-4-156}。
具体而言，先证明定理~\ref{thm:3-8-71}
的 $A_{n;K_1,K_2}$ 版本，
以证明定理~\ref{thm:7-4-154}、\ref{thm:7-4-156}
的 $A_{n;K_1,K_2}$ 版本。
关于 $A_{n;K_1,K_2}$ 双模、
$A_{n;K_1,K_2}$ 双模同态等概念的定义，
见下一小节。
把这些结果与若干同调代数结合，
即可证明定理~\ref{thm:7-4-154} 的
$A_\infty$ 版本，
从而证明定理~\ref{thm:3-8-71} 的
$A_\infty$ 版本。

为了得到定理~\ref{thm:7-4-156} 的
$A_\infty$ 版本，
还须证明由该定理的
$A_{n;K_1,K_2}$ 版本得到的同态彼此相容。
采用逐时积模空间的双参数版本，
即可证明这一点。
\end{remark}
